\section{Limitations and future work}
\label{sec:limitations}

\paragraph{Identification.} Volatility and regime risk premia are not identified
by a monthly baseload strip, and no attempt is made here to estimate them. All
reported values are zero-premium values in the sense of Section~\ref{sec:q1}.
A drift-type premium is additionally quadratically suppressed within this
specification, so even option data might identify it only weakly; an
intensity-type premium is the more promising target and is not implemented.

\paragraph{Parameter provenance.} Three items in Table~\ref{tab:provenance}
carry unresolved questions. The transition intercepts are moment-matching
reconstructions without standard errors. The transition slopes come from a
specification containing a ramp covariate that the production implementation
omits, with omitted-variable consequences of unknown sign. The transform scale
$s_P=282.48$ could not be verified because the preprocessing metadata defining it
was unavailable; it is a median absolute price over a training window spanning
years of substantial lira depreciation, so the price regime it represents may
differ materially from that at the valuation date. Each would be resolved by
access to the original estimation artefacts.

\paragraph{Specification scope.} The production residual is a single-factor
mean-reverting process, unbounded below, whereas realised prices face
administrative floors and ceilings. A companion specification developed alongside
this work adds a slow second factor and explicit price limits. It is not
integrated here because clipping the state distribution breaks the centring
identity of Proposition~\ref{prop:centering}: restoring $\E^{\Q}[P_t]=F(t)$ under
a clipped law requires redesigning the centring term rather than adding a bound
to the existing one. The trade-off is worth stating explicitly, because the
obvious alternative is worse. Multiplicative or log-type specifications respect
the lower bound automatically, but under Jensen's inequality their first moment
depends on the residual variance, so forward consistency can no longer be
imposed algebraically. Figure~\ref{fig:legacy} quantifies what that costs on
this data, and the forward inconsistency it shows is larger by orders of
magnitude, measured in \TRY, than the boundary violations the additive form
permits, whose simulated probability at the calibrated dispersion is zero to
four decimal places. That is the most substantive extension available and is
treated as the natural next study rather than a patch to this one. The dispersion
ratios in Table~\ref{tab:realised} indicate what a second factor would need to
supply. The price limit schedule used by the companion specification was inferred
from realised price behaviour rather than read from a primary regulatory
document, and is marked accordingly; it plays no role in the production results.

\paragraph{Evaluation.} The option-level evaluation still rests on a single
valuation date, because the residual parameters cannot be replayed at earlier
dates without look-ahead bias. Only the forward curve is evaluated across the
seven dates of Section~\ref{sec:multidate}, and the two exercises therefore
answer different questions. Seven dates six months apart also leave the delivery
windows partially overlapping and the long horizons confounded with calendar
month, so the horizon profile in particular should be read as descriptive.
Four of the seven strips were published one to four days before their valuation
date because the date fell on a non-business day; a strict same-day rule would
reduce the sample to three. Extending the panel as further quotations and
outturns accumulate, and re-estimating the residual process at each date so that
option prices could also be evaluated out of sample, are the natural
continuations.
